% GENERATED FILE. Edit canonical lessons, then run npm run generate:books.
% canonical-source-hash: fnv1a64:be7d17a6db502a4e
% canonical-lessons: ZH-C212-dayin, ZH-C212-fuyin, ZH-C212-chongzhi, ZH-C212-yuyue, ZH-C212-lifa

\chapter{Services in Town}
\label{ch:zh-services-in-town}
\hlchaptermodality{212}

\begin{chapteropening}
\textbf{By the end of this chapter:} \emph{I can print and photocopy things, top up my phone, make an appointment and get a haircut.}

Complete the last lesson of chapter 212: I can print and photocopy things, top up my phone, make an appointment and get a haircut.
\end{chapteropening}

\section[dǎyìn]{dǎyìn — to print out}

\label{lesson:ZH-C212-dayin}



\begin{quote}
Before the new one: say the Mandarin for map, then the Mandarin for everywhere.
\end{quote}

\subsection*{You'll want to know: dǎyìn}
\emph{dǎyìn} — \textquotedblleft{}to print out\textquotedblright{}.

\begin{quote}
\emph{wǒ xiǎng dǎyìn liǎng zhāng zhàopiàn} — I'd like to print two photos.
\end{quote}

Its characters come later; for now it is a word for the ear.

\begin{cousinweb}
\emph{dǎyìn} is built from \zh{打} 'to strike' and \emph{yìn} 'a seal', both inherited from Middle Chinese, \emph{taengX} and \emph{'jinH} (Guangyun, in Baxter's transcription), which glosses the second as a seal of office.
\end{cousinweb}

\begin{grammarlens}[title={What you've built}]
One more word for everyday actions.
\end{grammarlens}

\subsection*{Your turn}
\begin{itemize}
\raggedright
  \item \emph{Say it:} \emph{dǎyìn}
  \item \emph{Say it:} \emph{dǎyìn}, once more
  \item \emph{Say it:} \emph{wǒ xiǎng dǎyìn liǎng zhāng zhàopiàn} — I'd like to print two photos.
  \item \emph{Recall:} say the Mandarin for map, then the Mandarin for everywhere, then say \emph{dǎyìn} again
\end{itemize}

\begin{tcolorbox}[breakable,colback=teal!4,colframe=teal!35!black,title={Before you move on}]
What does dǎyìn mean? (\textquotedblleft{}To print out\textquotedblright{}.) Say it once more.
\end{tcolorbox}
\section[fùyìn]{fùyìn — to photocopy}

\label{lesson:ZH-C212-fuyin}



\begin{quote}
Before the new one: say the Mandarin for everywhere, then the Mandarin for to print out.
\end{quote}

\subsection*{You'll want to know: fùyìn}
\emph{fùyìn} — \textquotedblleft{}to photocopy\textquotedblright{}.

\begin{quote}
\emph{qǐng bāng wǒ fùyìn shēnfènzhèng} — Could you copy my identity card for me, please?
\end{quote}

Its characters come later; for now it is a word for the ear.

\begin{cousinweb}
\emph{fùyìn} is built from \emph{fù} 'double' and \emph{yìn} 'a seal, a print', both inherited from Middle Chinese, \emph{pjuwk} and \emph{'jinH} (Guangyun, in Baxter's transcription), which glosses the first as a lined garment.
\end{cousinweb}

\begin{grammarlens}[title={What you've built}]
One more word for everyday actions.
\end{grammarlens}

\subsection*{Your turn}
\begin{itemize}
\raggedright
  \item \emph{Say it:} \emph{fùyìn}
  \item \emph{Say it:} \emph{fùyìn}, once more
  \item \emph{Say it:} \emph{qǐng bāng wǒ fùyìn shēnfènzhèng} — Could you copy my identity card for me, please?
  \item \emph{Recall:} say the Mandarin for everywhere, then the Mandarin for to print out, then say \emph{fùyìn} again
\end{itemize}

\begin{tcolorbox}[breakable,colback=teal!4,colframe=teal!35!black,title={Before you move on}]
What does fùyìn mean? (\textquotedblleft{}To photocopy\textquotedblright{}.) Say it once more.
\end{tcolorbox}
\section[chōngzhí]{chōngzhí — to top up (a card or a phone)}

\label{lesson:ZH-C212-chongzhi}



\begin{quote}
Before the new one: say the Mandarin for to print out, then the Mandarin for to photocopy.
\end{quote}

\subsection*{You'll want to know: chōngzhí}
\emph{chōngzhí} — \textquotedblleft{}to top up (a card or a phone)\textquotedblright{}.

\begin{quote}
\emph{nǐ de shǒujī chōngzhí le ma} — Have you topped up your phone?
\end{quote}

Its characters come later; for now it is a word for the ear.

\begin{cousinweb}
\emph{chōngzhí} is built from \emph{chōng} 'to fill' and \emph{zhí} 'a value', both inherited from Middle Chinese, \emph{tsyhuwng} and \emph{driH} (Guangyun, in Baxter's transcription), which glosses the first 'full' and the second 'to hold' and 'to match'. The dictionary gives the second a departing tone, so the rising tone of zhí is not the regular outcome, and how it came is not settled.
\end{cousinweb}

\begin{grammarlens}[title={What you've built}]
One more word for everyday actions.
\end{grammarlens}

\subsection*{Your turn}
\begin{itemize}
\raggedright
  \item \emph{Say it:} \emph{chōngzhí}
  \item \emph{Say it:} \emph{chōngzhí}, once more
  \item \emph{Say it:} \emph{nǐ de shǒujī chōngzhí le ma} — Have you topped up your phone?
  \item \emph{Recall:} say the Mandarin for to print out, then the Mandarin for to photocopy, then say \emph{chōngzhí} again
\end{itemize}

\begin{tcolorbox}[breakable,colback=teal!4,colframe=teal!35!black,title={Before you move on}]
What does chōngzhí mean? (\textquotedblleft{}To top up (a card or a phone)\textquotedblright{}.) Say it once more.
\end{tcolorbox}
\section[yùyuē]{yùyuē — to make an appointment; to book}

\label{lesson:ZH-C212-yuyue}



\begin{quote}
Before the new one: say the Mandarin for to photocopy, then the Mandarin for to top up.
\end{quote}

\subsection*{You'll want to know: yùyuē}
\emph{yùyuē} — \textquotedblleft{}to make an appointment; to book\textquotedblright{}.

\begin{quote}
\emph{kěyǐ yùyuē ma} — Can I make an appointment?
\end{quote}

Its characters come later; for now it is a word for the ear.

\begin{cousinweb}
\emph{yùyuē} is built from \emph{yù} 'beforehand' and \emph{yuē} 'to bind, to agree', both inherited from Middle Chinese, \emph{yoH} and \emph{'jak} (Guangyun, in Baxter's transcription), which glosses the first 'before' and the second 'to tie up'.
\end{cousinweb}

\begin{grammarlens}[title={What you've built}]
One more word for everyday actions.
\end{grammarlens}

\subsection*{Your turn}
\begin{itemize}
\raggedright
  \item \emph{Say it:} \emph{yùyuē}
  \item \emph{Say it:} \emph{yùyuē}, once more
  \item \emph{Say it:} \emph{kěyǐ yùyuē ma} — Can I make an appointment?
  \item \emph{Recall:} say the Mandarin for to photocopy, then the Mandarin for to top up, then say \emph{yùyuē} again
\end{itemize}

\begin{tcolorbox}[breakable,colback=teal!4,colframe=teal!35!black,title={Before you move on}]
What does yùyuē mean? (\textquotedblleft{}To make an appointment; to book\textquotedblright{}.) Say it once more.
\end{tcolorbox}
\section[lǐfà]{lǐfà — to get a haircut}

\label{lesson:ZH-C212-lifa}



\begin{quote}
Before the new one: say the Mandarin for to top up, then the Mandarin for to make an appointment.
\end{quote}

\subsection*{You'll want to know: lǐfà}
\emph{lǐfà} — \textquotedblleft{}to get a haircut\textquotedblright{}.

\begin{quote}
\emph{wǒ yùyuēle míngtiān qù lǐfà} — I've booked to get my hair cut tomorrow.
\end{quote}

Its characters come later; for now it is a word for the ear.

\begin{cousinweb}
\emph{lǐfà} is built from \emph{lǐ} 'to put in order' and \emph{fà} 'the hair of the head', both inherited from Middle Chinese, \emph{liX} and \emph{pjot} (Guangyun, in Baxter's transcription), which glosses the first 'to work jade'.
\end{cousinweb}

\begin{grammarlens}[title={What you've built}]
One more word for everyday actions.
\end{grammarlens}

\subsection*{Your turn}
\begin{itemize}
\raggedright
  \item \emph{Say it:} \emph{lǐfà}
  \item \emph{Say it:} \emph{lǐfà}, once more
  \item \emph{Say it:} \emph{wǒ yùyuēle míngtiān qù lǐfà} — I've booked to get my hair cut tomorrow.
  \item \emph{Recall:} say the Mandarin for to top up, then the Mandarin for to make an appointment, then say \emph{lǐfà} again
\end{itemize}

\begin{tcolorbox}[breakable,colback=teal!4,colframe=teal!35!black,title={Before you move on}]
What does lǐfà mean? (\textquotedblleft{}To get a haircut\textquotedblright{}.) Say it once more.
\end{tcolorbox}
